\section{Estado del Arte}

\subsection{Estado de madurez del proceso y evolución esperada}

En la actualidad, el entorno operativo de las instituciones deportivas en Argentina se caracteriza por un grado de madurez tecnológica incipiente y fuertemente fragmentado. La gestión societaria de la mayoría de los clubes se basa de manera predominante en herramientas administrativas reactivas, procesos organizativos manuales y la emisión de credenciales físicas altamente vulnerables a la falsificación. Este paradigma analógico o de baja digitalización genera ineficiencias críticas que desfinancian de manera crónica a los clubes: una alta tasa de morosidad y deserción societaria silenciosa, fraude perimetral mediante la duplicación de accesos los días de partido, y un desaprovechamiento sustancial de los activos institucionales.

Frente a este \textit{status quo}, la evolución esperada se materializa en el ecosistema de \textbf{SocioUnido}, una plataforma \textit{Software como Servicio (SaaS)} de tipo \textit{de empresa a empresa (E2E)}. La transición operativa propone abandonar la administración estática para adoptar un modelo predictivo, proactivo y escalable. Esta evolución se fundamenta en componentes de alta ingeniería de software: la integración de un motor analítico basado en \textit{aprendizaje automático (Machine Learning)} para la detección temprana de anomalías en la asistencia y riesgos de baja; la implementación de un asistente transaccional con procesamiento de lenguaje natural (PLN) vía WhatsApp para lograr una fricción cero en las transacciones; y la proyección de un acceso inteligente (\textit{Smart Access}) amparado en algoritmos criptográficos TOTP, permitiendo validaciones 100\% fuera de línea (\textit{offline}) para tolerar el colapso sistemático de las redes de telefonía móvil.

\subsection{Definición arquitectónica y marco de trabajo}

El proyecto ha superado su fase de concepción inicial, encontrándose en un estado de madurez donde su estructura técnica y directrices operativas están plenamente establecidas. La arquitectura del ecosistema, basada en microservicios, ya se encuentra completamente diseñada y documentada (su modelado detallado puede consultarse en el \textbf{Anexo B} del presente anteproyecto).

En consonancia con los estándares contemporáneos de la industria del software, el desarrollo de la plataforma se rige bajo la metodología ágil iterativa e incremental \textit{Scrum}. El trabajo se planifica y ejecuta mediante \textit{Sprints} semanales de estabilización, garantizando entregas continuas de valor, adaptabilidad ante nuevos requerimientos y un control de calidad riguroso sobre cada incremento del producto.

\newpage

\subsection{Ecosistema de herramientas y gestión documental}

Para sostener el más alto nivel académico, asegurar la trazabilidad del proceso y facilitar la colaboración, la totalidad de la gestión del código fuente, el control de versiones y la planificación del Trabajo Profesional se ha centralizado en \textbf{GitHub}. Todo el entorno opera bajo el paraguas de una misma ``Organización de GitHub'', lo cual permite articular de forma nativa los repositorios, las historias de usuario, los flujos de integración y entrega continua (CI/CD) y el seguimiento de incidencias.

Este ecosistema organizativo se encuentra compartimentado estratégicamente de la siguiente manera:

\begin{itemize}
    \item \textbf{Repositorios de desarrollo técnico:} Espacios independientes y aislados para cada microservicio y para las interfaces de usuario (tanto el panel administrativo como la aplicación web progresiva). En cada repositorio reside la implementación lógica, las rutinas de pruebas (\textit{testing}) y la documentación técnica específica del componente.
    \item \textbf{Repositorios de documentación interactiva:} Se emplea la herramienta \textit{Just the Docs} para automatizar la generación de páginas web estáticas y amigables. Estas páginas documentan de forma accesible procesos iterativos clave: la bitácora del proyecto, la gestión \textit{Scrum} (ceremonias, hoja de ruta o \textit{Roadmap}, y métricas de rendimiento), y el registro analítico de todas las ideas de proyecto exploradas durante la etapa formativa.
    \item \textbf{Repositorio de entregables académicos:} Un entorno centralizado de almacenamiento estático para toda la documentación formal e institucional. Aquí se resguardan documentos como este Anteproyecto, los archivos multimedia de las entregas intermedias, y cualquier guía o análisis complementario requerido por la cátedra a medida que avanza el ciclo de desarrollo. Además de por supuesto, la documentación final de todo el Trabajo Profesional que se entregará al concluir el proyecto.
    \item \textbf{Centros de monitoreo integrados (GitHub Projects):} Para la coordinación operativa, se han desplegado dos proyectos independientes que funcionan como centros de comando:
    \begin{itemize}
        \item \textit{Centro de monitoreo de tareas:} Tablero dedicado exclusivamente a la gestión de funcionalidades, historias de usuario y priorización del \textit{Product Backlog}.
        \item \textit{Centro de monitoreo de incidencias (\textit{Bugs}):} Entorno específico para el rastreo, asignación y resolución de errores identificados durante las fases de validación y estabilización del software.
    \end{itemize}
\end{itemize}
